\chapter{Introduction}

In the vastness of entomology, the study of insect flight mechanics involving a beautifully orchestrated dance between muscles, exoskeleton and flexible wings has been a topic of interest in bio-mechanics research. Over the years, several research works have been trying to understand and learn from seemingly complex manoeuvres that defy the conventional notions of efficiency and performance carried out by insects\cite{Farisenkov2022,range}. A set of muscles, exoskeletal structures, and wings known as the flight motor\cite{Pons_2023} play vital role in achieving this performance. Using different mechanisms across different species \cite{PRINGLE1968163}, the flight motor generates complex flapping wing beat kinematics \cite{kinematics1,kinematics2} which generate lift. The details of these kinematics are not fully understood\cite{Pons_2023}. Such interplay of active (due to muscle action) and passive elements (due to fluid-structure interaction) of the insect flight mechanics is particularly captivating of the seemingly common \textit{Drosophila melanogaster (\textit{Fruit Fly})}\cite{Pons_2023}. Insects have a set of flight steering muscles to control their wing beat kinematics\cite{Balint}, but they also probably rely on passive fluid-structure interaction to generate certain kinematic features \cite{PhysRevE.92.022712}. In particular, the wing pitch (or, incidence) degree of freedom (DOF) is thought to be passive in fruit flies \cite{PhysRevE.92.022712} – but only on the basis of aerodynamic evidence from quasi-steady aerodynamic models. To confirm whether the pitch DOF is passive, we need higher-fidelity aerodynamic data, e.g. from CFD. In this report, an investigation carried out using CFD on the unsteady aerodynamics of a rigid fruit fly wing executing hovering motion is presented. The work was primarily aimed at determining the degree of passiveness of different degrees of freedom involved in the fruit fly wing beat cycle.

\section{General Scope}

Among different domains of applications, optimal control systems help achieve efficiency and performance in operation. Minimizing the number of active controls may reduce the chances of error and inefficiency in such systems. With this background, we have been looking into nature and studying different passive motions in it. In the present work, we focus on studying the complex wing beat motion of a fruit fly. The fruit fly wing beat is characterized to have three degrees of freedom namely pitch, stroke, and elevation \cite{kinematics1}. The general scope of this study is to understand the complex wing beat motion and analyze the passiveness of different degrees of freedom.
\\
\\
We can approach the study in two ways. By an inverse approach where we iterate on stiffness damping properties of the fruit fly wing model until we achieve convergence with the observed kinematics. The second approach is where we use CFD to simulate the wing kinematics and find the stiffness and damping properties in different degrees of motion. In this work, we discuss and build on the CFD approach. The CFD approach is carried out using two types of simulations. First, the approach which is discussed in detail in this report is by considering a rigid fruit fly wing that follows the kinematics.
The second approach where we do a fully coupled CFD simulation with flexible wings is considered as a future work. By simulating the aerodynamics of rigid fruit fly wings, we aim to identify the passiveness of pitch motion in particular.


\section{Methodology}
The kinematics of the fruit-fly wing motion is taken from literature \cite{kinematics1,kinematics2}. Fruit fly wing geometry is taken from literature\cite{kinematics1}. The motion kinematics were defined using Fourier series approximations. The entire wing beat cycle was discretized and Euler angles for the wing for different instances along the discretized wing beat cycle are calculated from the Fourier approximations. This data is used to simulate and analyze the wing beat motion in CFD software. STAR-CCM+ software and cluster computing platforms were used for the simulations. 

\section{Expected Outcomes}

Insights on different patterns of motions involved as well as concluding on the degree of passiveness of different degrees of freedom are the primary expected outcomes of this project. Also, We expect the simulations to aid in setting up a fully coupled simulation of the fruit fly wing and also to develop an accurate model of the complex fruit fly wing motion. The method used in this approach can be utilized to simulate and study similar unsteady flapping motions.


